% ===================================================================
%  BAGAN No. 3 — Masa Kanak-kanak dan Masa Muda.
%
%  Traduction, faite en 2026, en indonesien d'aujourd'hui, sur l'ido
%  de texto/io. Regles de la colonne : voir 10-bagan-01.tex. Deux
%  scenes, deux numerotations : les cles portent c1 et c2. Le
%  fac-simile ne met pas « julio » en gras la ou il gras « junio » et
%  « agosto » ; la correction est declaree dans gravuri/korekti.json
%  et la colonne ecrit « (Juli atau », qui est la forme que cette
%  correction attend.
%
%  LA PREMIERE COLONNE DU RELEVE OU CE SONT LES SAISONS QUI MANQUENT.
%  Toutes les colonnes precedentes butaient sur des CHOSES — le marron
%  chaud, le fleau, le tilleul, la bruyere. Ici c'est le cadre du
%  tableau lui-meme : l'Indonesie est equatoriale et n'a pas quatre
%  saisons mais deux, musim hujan et musim kemarau, la saison des
%  pluies et la saison seche. « musim semi » (le printemps) et
%  « musim gugur » (l'automne) existent bel et bien en indonesien —
%  ce sont des mots d'ecole, forges pour la lecon de geographie, que
%  tout le monde comprend et que personne n'a vecu.
%
%  ON LES ECRIT SANS DETOUR, ET C'EST LA MEME REGLE QUE POUR LE VIN
%  ET LE JEU DE BARRES. Ecrire « musim hujan » a la place du
%  printemps donnerait un tableau ou les hirondelles reviennent en
%  mars sous la mousson : ce serait transposer la chose, pas traduire
%  la langue. La colonne garde donc mars, avril, mai, le printemps et
%  l'ete du fac-simile, et laisse le lecteur indonesien devant un
%  calendrier qui n'est pas le sien — exactement comme le lecteur
%  francais de 1926 etait devant un lion et une girafe.
%
%  ET LA MENAGERIE RENVERSE LE RAPPORT. Au bloc t03-c2-15-1, l'ido
%  montre « la sovaja bestii » comme une curiosite de foire :
%  harimau le tigre (51), gajah l'elephant (54). Ce sont deux
%  animaux de Sumatra et de Java, nommes par des mots malais de fond,
%  et l'ecolier indonesien de 2026 les connait mieux que l'ecolier
%  francais de 1926. Seule la girafe (53) reste etrangere — jerapah,
%  emprunte a l'arabe zarafa, comme le francais. Ce tableau est le
%  premier ou l'exotisme change de camp.
%
%  singa, LE LION (52), MERITE SA PHRASE. Le mot vient du sanskrit
%  simha, il est arrive avec l'hindouisme il y a quinze siecles, et il
%  a donne son nom a Singapura, « la ville du lion » — dans une
%  region ou le lion n'a jamais vecu. Le mot a voyage sans la chose,
%  et il est aujourd'hui plus indonesien que beaucoup de mots pour des
%  betes qu'on voit. C'est l'exact contraire du glacier ourdou du
%  tableau 9 : la, la chose sans le mot ; ici, le mot sans la chose.
%
%  LES VETEMENTS SONT LE SEUL ENDROIT OU LA COUCHE NEERLANDAISE
%  REVIENT EN FORCE : rompi le gilet, dasi la cravate, bretel les
%  bretelles, manset les manchettes, korset le corset, jas le
%  redingote. Le tableau 2 avait montre le corps entierement malais ;
%  le tableau 3 montre ce qu'on met dessus entierement neerlandais.
%  Le partage n'est pas entre l'ancien et le moderne, il est entre ce
%  qu'on a et ce qu'on achete.
% ===================================================================

%%K t03-apar-1 apar
\VUsaut{3.0mm}
\VUtitre{15.0pt}{60}{BAGAN No. 3}
\VUsaut{1.6mm}
\VUtitre{11.4pt}{0}{Masa Kanak-kanak dan Masa Muda.}
\VUsaut{3.4mm}

%%K t03-apar-2 apar
(Bagan ke-3 dan ke-4 disusun sedemikian rupa sehingga keduanya
menyajikan, dalam empat \VUgras{adegan keluarga}, pengertian-pengertian
pokok tentang \VUgras{cuaca}, \VUgras{umur}, \VUgras{keluarga},
\VUgras{pakaian} dan \VUgras{keadaan}.)

\VUsaut{4.0mm}
%%K t03-tit-1 sub
\VUcentre{12.0pt}{0}{\textit{Adegan Pertama.}}
\VUsaut{3.0mm}

%%K t03-tit-2 sub
\VUpk{10.2pt}{40}{Upacara pembaptisan di sebuah desa.}
\VUsaut{2.4mm}

%%K t03-tit-3 sub
\VUpk{10.2pt}{40}{Musim semi.}
\VUsaut{3.0mm}

%%K t03-c1-01-1 p
1. --- Kita berada pada \VUgras{musim semi}, di bulan \VUgras{Maret},
\VUgras{April} atau \VUgras{Mei}, pada hari-hari \VUgras{minggu},
\VUgras{Senin}, \VUgras{Selasa} atau \VUgras{Rabu}. Waktu menunjukkan
\VUgras{pukul sepuluh} pagi.

%%K t03-c1-02-1 p
2. --- \VUgras{Cuaca} indah, \VUgras{udara} bersih. Matahari
bersinar. Beberapa \VUgras{awan kecil} \textsuperscript{(2)} naik di
\VUgras{langit} \textsuperscript{(1)} yang biru.

%%K t03-c1-03-1 p
3. --- \VUgras{Pohon-pohon} hampir belum berdaun. \VUgras{Pohon
poplar} \textsuperscript{(3)} yang langsing dan \VUgras{pohon platan}
\textsuperscript{(17)} yang tinggi sampai sekarang baru mempunyai
kuncup.

%%K t03-c1-04-1 p
4. --- \VUgras{Burung-burung layang-layang} \textsuperscript{(4)}
kembali dan beterbangan dengan cicit yang gembira di sekitar
\VUgras{menara runcing} \textsuperscript{(7)} pada \VUgras{menara
lonceng} \textsuperscript{(5)} yang memahkotai \VUgras{gereja}
\textsuperscript{(6)} desa itu.

%%K t03-tit-4 sub
\VUsaut{3.4mm}
\VUpk{10.2pt}{40}{Gereja.}
\VUsaut{3.0mm}

%%K t03-c1-05-1 p
5. --- Gereja itu sangat sederhana, dengan \VUgras{gerbangnya}
\textsuperscript{(13)} yang besar dan \VUgras{jamnya}
\textsuperscript{(8)} yang besar. \VUgras{Jarum kecil}
\textsuperscript{(9)} ada pada angka X ; ia menunjukkan \VUgras{jam.}
\VUgras{Jarum besar} \textsuperscript{(10)} ada pada XII (tengah
lingkaran) ; ia menunjukkan \VUgras{menit.}

%%K t03-c1-06-1 p
6. --- Di menara lonceng, \VUgras{lonceng-lonceng}
\textsuperscript{(11)} berbunyi dengan riang. Di puncak atap berdiri
sebuah \VUgras{salib} \textsuperscript{(12)} kecil dari batu.

%%K t03-c1-07-1 p
7. --- Gereja itu bukan hanya mempunyai \VUgras{gerbang}
\textsuperscript{(13)} besar ; ia juga mempunyai pintu kecil di
samping. Melalui pintu kecil itu \VUgras{pastor} \textsuperscript{(15)},
berpakaian \VUgras{jubahnya}, keluar dari \VUgras{ruang sakristi}
\textsuperscript{(14)} dan pergi ke \VUgras{pastoran}
\textsuperscript{(19)}.

%%K t03-c1-08-1 p
8. --- Di dekatnya, inilah \VUgras{penjual susu}
\textsuperscript{(24)} dengan \VUgras{keledainya} \textsuperscript{(23)}.
Ia kembali dari kota, tempat ia membawa susu yang baik untuk
anak-anak.

%%K t03-tit-5 sub
\VUsaut{4.0mm}
\VUpk{10.2pt}{40}{Kebun buah.}
\VUsaut{3.4mm}

%%K t03-c1-09-1 p
9. --- Tuan Aliboron (si keledai) sangat puas dengan perjalanan pagi
itu. Ia menghirup dengan nikmat udara yang diharumkan oleh wangi
\VUgras{bunga-bunga} \textsuperscript{(20)} yang menutupi \VUgras{pohon
buah} \textsuperscript{(18)} di \VUgras{kebun buah} sebelah. Merah muda
dan putih benar \VUgras{pohon-pohon persik} \textsuperscript{(18)} dan
\VUgras{pohon-pohon aprikot} \textsuperscript{(19)} itu, dan mereka
menjanjikan panen yang baik.

%%K t03-c1-10-1 p
10. --- Sementara itu \VUgras{lebah-lebah} \textsuperscript{(22)} kecil
dari \VUgras{sarangnya} \textsuperscript{(21)} pergi ke sana mencuri
mengumpulkan \VUgras{madu} yang manis sambil berdengung.

%%K t03-c1-11-1 p
11. --- Tetapi apa yang terjadi? Inilah anak-anak desa berlarian dan
membuat \VUgras{angsa-angsa} \textsuperscript{(25)} lari ketakutan.
Itulah keluarnya rombongan dari gereja sesudah \VUgras{pembaptisan}
anak lelaki notaris.

%%K t03-c1-12-1 p
12. --- Si kecil Iakobus, di dalam \VUgras{ayunan bayinya}
\textsuperscript{(26)}, terbangun oleh bunyi lonceng yang riang, dan
kakak perempuannya, Magdalena, yang juga ingin melihat rombongan
pembaptisan, meminta ibunya datang menyisir rambutnya dan
memakaikan pakaiannya di ambang pintu rumah.

%%K t03-c1-13-1 p
13. --- Ibunya menyetujui. Ia memakai \VUgras{selendangnya}
\textsuperscript{(28)}, \VUgras{baju atasannya} \textsuperscript{(29)}
dan \VUgras{celemeknya} \textsuperscript{(30)}, lalu duduk di sebuah
bangku kecil di depan pintu. Magdalena baru memakai \VUgras{rok
dalamnya} \textsuperscript{(31)} dan \VUgras{korsetnya}
\textsuperscript{(32)}.

%%K t03-c1-14-1 p
14. --- Betapa bahagianya ia! Ia melupakan bunga-bunga
kesayangannya : \VUgras{bunga anyelirnya} \textsuperscript{(34)}, tempat
\VUgras{kupu-kupu} \textsuperscript{(35)} hinggap, \VUgras{bunga
tulipnya} \textsuperscript{(36)}, \VUgras{bunga bakung Meinya}
\textsuperscript{(37)} dalam \VUgras{pot} \textsuperscript{(38)} di atas
\VUgras{tembok} \textsuperscript{(33)}, dan \VUgras{bunga mawarnya}
\textsuperscript{(39)} yang membentuk pagar begitu indah. Ia
memperhatikan pakaian mewah para nyonya dalam iringan itu.

%%K t03-c1-15-1 p
15. --- Anak-anak desa tidak begitu memperhatikan pakaian. Mereka
berdesak-desakan dan saling dorong untuk mendapat \VUgras{gula-gula}
\textsuperscript{(55)} atau \VUgras{uang logam} \textsuperscript{(52)}.
Salah seorang dari mereka menangis \textsuperscript{(40)}.

%%K t03-c1-16-1 p
16. --- Yang lain menginjak kaki \VUgras{anjing} \textsuperscript{(42)}
yang lari sambil menyalak dan membuat \VUgras{induk ayam}
\textsuperscript{(43)} beserta \VUgras{anak-anak ayamnya}
\textsuperscript{(44)} lari pula.

%%K t03-tit-6 sub
\VUsaut{5.0mm}
\VUpk{10.2pt}{40}{Iringan pembaptisan.}
\VUsaut{4.0mm}

%%K t03-c1-17-1 p
17. --- Di depan iringan berjalan \VUgras{inang pengasuh}
\textsuperscript{(45)}. Ia memakai \VUgras{tudung kepala}
\textsuperscript{(46)} yang indah dengan pita-pita panjang. Ia
menggendong \VUgras{bayi yang baru lahir} \textsuperscript{(47)}, yang
seluruhnya berpakaian \VUgras{renda} \textsuperscript{(48)}.

%%K t03-c1-18-1 p
18. --- Di belakang inang pengasuh datang \VUgras{bapak baptis}
\textsuperscript{(49)} dan \VUgras{ibu baptis} \textsuperscript{(53)}.
Merekalah yang membagikan gula-gula dan uang logam. Bapak baptis
memakai \VUgras{sarung tangan} \textsuperscript{(51)} dan \VUgras{topi
tinggi} \textsuperscript{(50)}. Ibu baptis memegang \VUgras{karangan
bunga} \textsuperscript{(54)} yang indah di tangannya.

%%K t03-c1-19-1 p
19. --- Di belakang mereka kita melihat \VUgras{paman}
\textsuperscript{(56)} dan \VUgras{bibi} \textsuperscript{(58)} bayi
itu. Bibi memakai \VUgras{topi} \textsuperscript{(59)} yang elegan,
paman memakai \VUgras{jas panjang} \textsuperscript{(57)} hitam dan
\VUgras{rompi} putih. Lalu inilah \VUgras{sepupu laki-laki}
\textsuperscript{(60)} dan \VUgras{sepupu perempuan}
\textsuperscript{(61)}. Yang terakhir ini memegang tangan
\VUgras{kakak perempuan} \textsuperscript{(62)} si bayi, si kecil
Berta.

%%K t03-c1-20-1 p
20. --- \VUgras{Saudara laki-laki} \textsuperscript{(63)} Berta yang
lain berjalan di belakang. Ia memberikan tangannya kepada seorang
\VUgras{kerabat perempuan} \textsuperscript{(64)}. \VUgras{Teman-teman}
\textsuperscript{(65)} keluarga itu datang kemudian.

%%K t03-tit-7 sub
\VUsaut{4.6mm}
\VUpk{10.2pt}{40}{Para penonton.}
\VUsaut{3.4mm}

%%K t03-c1-21-1 p
21. --- Di dekat rombongan, inilah seorang \VUgras{pengemis}
\textsuperscript{(66)} yang bertumpu pada \VUgras{tongkat ketiaknya}
\textsuperscript{(67)}. Ia meminta sedekah.

%%K t03-c1-22-1 p
22. --- Di sisi lain ada seorang anak laki-laki kecil yang sangat
\VUgras{sopan} \textsuperscript{(68)}. Ia memberi salam dan mengucapkan
« \VUgras{selamat pagi} » kepada orang-orang yang dikenalnya.

%%K t03-c1-23-1 p
23. --- Orang-orang desa keluar dari rumah mereka untuk melihat
iringan pembaptisan. Kita melihat seorang \VUgras{orang desa}
\textsuperscript{(69)} dengan \VUgras{peci katunnya} dan di sebelahnya
seorang \VUgras{perempuan desa} \textsuperscript{(70)}. Kemudian
seorang \VUgras{perempuan tua} \textsuperscript{(71)} yang bertumpu
pada \VUgras{tongkatnya} \textsuperscript{(72)}. Akhirnya
\VUgras{tukang giling} \textsuperscript{(73)} dengan \VUgras{topi
lakennya} \textsuperscript{(74)} yang besar, dan seorang perempuan desa
muda beralas kaki \VUgras{terompah kayu} \textsuperscript{(75)}, yang
sedang mengajari anaknya berjalan.

%%K t03-c1-24-1 p
24. --- Adegan itu sangat menyenangkan.

\VUsaut{4.0mm}
%%K t03-tit-8 sub
\VUcentre{12.0pt}{0}{\textit{Adegan Kedua.}}
\VUsaut{3.0mm}

%%K t03-tit-9 sub
\VUpk{10.2pt}{40}{Pesta rakyat.}
\VUsaut{2.4mm}

%%K t03-tit-10 sub
\VUpk{10.2pt}{40}{Permainan air dan lomba dayung.}
\VUsaut{3.0mm}

%%K t03-c2-01-1 p
1. --- \VUgras{Musim panas} telah tiba. Matahari yang sangat terik di
bulan \VUgras{Juni} (Juli atau \VUgras{Agustus}) mendorong kaum muda
untuk berenang dan berperahu.

%%K t03-c2-02-1 p
2. --- Di tepi kanan sungai, para pendayung menanggalkan pakaian
untuk ikut serta dalam permainan air dan lomba dayung.

%%K t03-c2-03-1 p
3. --- Salah seorang dari mereka melepas \VUgras{sepatunya}
\textsuperscript{(1)} ; ia membuka talinya. Kedua kawannya sudah siap.
Sekarang mereka hanya memakai \VUgras{baju rajut}
\textsuperscript{(2)} dan \VUgras{celana renang} \textsuperscript{(3)}.
Mereka berbincang sambil menunggu teman-teman mereka. Mereka
berbicara tentang pasar malam dan pesta yang berlangsung di seberang.

%%K t03-c2-04-1 p
4. --- Di tanah, di dekat mereka, tergeletak \VUgras{pakaian}
kawan-kawan mereka : sepasang \VUgras{sepatu bot} \textsuperscript{(4)},
\VUgras{sepatu} \textsuperscript{(5)}, \VUgras{kaus kaki}
\textsuperscript{(6)}, topi \VUgras{laken} \textsuperscript{(7)} dan
topi \VUgras{jerami} \textsuperscript{(8)}, serta \VUgras{dasi}
\textsuperscript{(9)}.

%%K t03-c2-05-1 p
5. --- Salah seorang pemuda melipat \VUgras{setelannya}
\textsuperscript{(10)} dengan hati-hati. Ia meletakkannya di atas kain,
yang sebentar lagi akan dipakai tetangganya untuk membungkus kedua
\VUgras{pakaian} itu.

%%K t03-c2-06-1 p
6. --- Anak keenam terlambat. Ia masih memakai \VUgras{topi petnya}
\textsuperscript{(11)} di kepala. Ia belum melepas \VUgras{kemejanya}
\textsuperscript{(12)}, \VUgras{kerahnya} \textsuperscript{(13)},
maupun \VUgras{mansetnya} \textsuperscript{(14)}. Ia memegang
\VUgras{rompinya} \textsuperscript{(17)} di tangan dan masih memakai
\VUgras{celana panjangnya} \textsuperscript{(16)} dan
\VUgras{bretelnya} \textsuperscript{(15)}. Tentu ia tidak akan siap
untuk perlombaan berikutnya.

%%K t03-c2-07-1 p
7. --- Di dekat para pemuda, sebuah jalan setapak yang di kedua
sisinya banyak \VUgras{tumbuhan duri} \textsuperscript{(18)} menuju
tepi sungai, tempat bapak Iakobus menyiapkan di antara
\VUgras{buluh-buluh} \textsuperscript{(19)} \VUgras{kembang api}
\textsuperscript{(20)} yang akan dinyalakan malam nanti.

%%K t03-tit-11 sub
\VUsaut{4.4mm}
\VUpk{10.2pt}{40}{Perlombaan.}
\VUsaut{3.4mm}

%%K t03-c2-08-1 p
8. --- Di sungai, beberapa nyonya, berlindung di bawah payung mereka,
berjalan-jalan dengan \VUgras{perahu} \textsuperscript{(21)}. Mereka
mengikuti perlombaan yang sangat menarik itu. Di setiap
\VUgras{perahu dayung} \textsuperscript{(22)}, empat
\VUgras{pendayung} \textsuperscript{(24)} mendayung sekuat tenaga,
sementara \VUgras{juru mudi} \textsuperscript{(26)} memegang
\VUgras{kemudi} \textsuperscript{(25)} dan mengarahkan perahu dayung
yang rapuh itu. \VUgras{Dayung-dayung} \textsuperscript{(23)} memukul
air dengan berirama dan perahu-perahu ringan itu meluncur dengan
kecepatan yang memusingkan.

%%K t03-c2-09-1 p
9. --- Beberapa pemuda bertanding, di dekat tiang jembatan,
memperebutkan hadiah \VUgras{tiang haluan} \textsuperscript{(28)}.
Salah seorang berjalan dengan hati-hati di atas tiang yang lentur dan
licin itu untuk meraih \VUgras{bendera} \textsuperscript{(29)} kecil
yang dipasang di ujungnya. \VUgras{Peserta} \textsuperscript{(30)}
lawannya yang sial jatuh ke air. Ia berenang untuk mencapai daratan.

%%K t03-c2-10-1 p
10. --- Pemuda-pemuda yang lebih tua \VUgras{bertanding perahu}
\textsuperscript{(32)} di dekat \VUgras{dermaga} \textsuperscript{(33)}.
Semua permainan itu disaksikan oleh \VUgras{orang banyak}
\textsuperscript{(31)} yang bertumpu pada \VUgras{tembok pembatas}
\VUgras{jembatan} \textsuperscript{(27)} dan berkerumun di
\VUgras{tepi} sungai. Namun beberapa penonton berbalik untuk
mengikuti permainan \VUgras{tiang hadiah} \textsuperscript{(45)} atau
untuk mengagumi \VUgras{iringan wali kota} yang datang untuk
\VUgras{pembagian} hadiah.

%%K t03-c2-11-1 p
11. --- Pertama-tama, inilah beberapa \VUgras{anak nakal}
\textsuperscript{(35)} yang mendahului \VUgras{para pemain musik}
\textsuperscript{(36)} ; kemudian datang \VUgras{wali kota}
\textsuperscript{(37)}, para wakilnya dan \VUgras{dewan kota} kota
kecil itu. Akhirnya berjalan \VUgras{para pemadam kebakaran}
\textsuperscript{(38)}.

%%K t03-tit-12 sub
\VUsaut{5.0mm}
\VUpk{10.2pt}{40}{Pasar malam.}
\VUsaut{3.6mm}

%%K t03-c2-12-1 p
12. --- Banyak sekali orang di \VUgras{alun-alun pasar malam}
\textsuperscript{(39)}. Orang datang melihat bermacam-macam
\VUgras{kios} : \VUgras{komidi putar} \textsuperscript{(40)} berkuda
kayu, \VUgras{sirkus} \textsuperscript{(41)} yang pertunjukannya sudah
dimulai, \VUgras{pesta dansa umum} \textsuperscript{(42)}, tempat
pemuda dan pemudi kota menikmati kesenangan \VUgras{menari.}

%%K t03-c2-13-1 p
13. --- Di ujung sana kita melihat \VUgras{arena}
\textsuperscript{(44)}, tempat \VUgras{matador} Spanyol yang berani
bertarung melawan \VUgras{banteng} \textsuperscript{(45)} yang
mengamuk ; kemudian \VUgras{luncuran rel} \textsuperscript{(49)}, dan
\VUgras{tempat menembak} \textsuperscript{(46)}, tempat orang berlatih
memakai \VUgras{senapan} dan menembak ke \VUgras{sasaran.}

%%K t03-c2-14-1 p
14. --- Di dekatnya, inilah \VUgras{teater pasar malam}
\textsuperscript{(47)}, tempat dimainkan \VUgras{komedi} dan
\VUgras{drama.} Sangat dekat dari situ ada kios \VUgras{raksasa}
\textsuperscript{(48)} dan \VUgras{orang kerdil.} \VUgras{Tinggi} yang
pertama dua meter lebih lima belas sentimeter ; yang kedua hampir
tidak lebih besar daripada seorang anak.

%%K t03-c2-15-1 p
15. --- Akhirnya, inilah \VUgras{kebun binatang keliling}
\textsuperscript{(50)} yang besar dengan \VUgras{binatang-binatang
buasnya}, \VUgras{harimaunya} \textsuperscript{(51)} dengan
\VUgras{cakar} yang kuat, \VUgras{singanya} \textsuperscript{(52)}
dengan \VUgras{surai} yang lebat, kedua \VUgras{jerapahnya}
\textsuperscript{(53)} dan \VUgras{gajahnya} \textsuperscript{(54)}
yang begitu cekatan memakai \VUgras{belalainya.}

%%K t03-c2-16-1 p
16. --- \VUgras{Penjinak binatang} \textsuperscript{(55)} yang melatih
binatang-binatang itu berdiri di pintu kios.

\VUsaut{6.0mm}
\VUfilet{22mm}
